\title{คาบเรียนที่ 8: การประยุกต์ของอนุพันธ์ (2/2)}
\frame{\titlepage}

\begin{frame}
ในคาบเรียนนี้ เราจะศึกษาบทประยุกต์ของอนุพันธ์ต่อ
\tableofcontents
\end{frame}


%%%%%%%%%%%%%%
\section[การกระจัด ความเร็ว และความเร่ง]{การกระจัด ความเร็วและความเร่งในเชิงอนุพันธ์}
\frame{\sectionpage}
%%%%%%%%%%%%%%

\begin{frame}
\addfig{0.8}{diff_applied_sva_intro}{การกระจัด ความเร็ว และความเร่งของรถยนต์จะแตกต่างไปในแต่ละเวลา \href{https://www.pexels.com/photo/red-racing-car-on-race-track-during-daytime-158971/}{(ที่มา)}}{}
\end{frame}

\begin{frame}
จากเรื่องของการเคลื่อนที่แบบโปรเจคไทล์ เราได้ทราบว่า ถ้าวัตถุมีความเร่งเป็นค่าคงที่เท่ากับ $a$ แล้ว 
\begin{itemize}
\item ความเร็วของวัตถุ ณ เวลา $t$ จะมีค่าเท่ากับ
\[v = at + v_0\]
เมื่อ $v_0$ คือความเร็วต้น
\item ตำแหน่งของวัตถุ ณ เวลา $t$ จะมีค่าเท่ากับ
\[s = \frac{1}{2}at^2 + v_0t + s_0\]
เมื่อ $s_0$ คือตำแหน่งเริ่มต้น
\end{itemize}
ในที่นี้ เราจะแสดงให้เห็นถึงที่มาของสูตรเหล่านี้
\end{frame}

\begin{frame}
สังเกตว่า ถ้าความเร่งเป็นค่าคงที่ แล้วความเร็วจะเป็นกราฟเส้นตรง และการกระจัดจะเป็นกราฟพาราโบลา
\addfig{0.9}{diff_applied_sva}{เปรียบเทียบกราฟ $s(t), v(t)$ และ $a(t)$}{}
\end{frame}

\begin{frame}[t]
\frametitle{ความสัมพันธ์ระหว่างการกระจัด $s$ และความเร็ว $v$}
สมมติว่าวัตถุเคลื่อนที่ตามแนวเส้นตรง \pa โดยมีสมการเคลื่อนที่ $s = s(t)$ เมื่อ $t$ แทนเวลา และ $s$ แทนระยะทาง\pa ความเร็วเฉลี่ยของวัตถุจาก $t = t_1$ ถึง $t = t_2$ \pa คือ 
\begin{slide}
\[ \frac{s(t_2) - s(t_1)}{t_2 - t_1}\] \pa

ถ้าเราหาลิมิตเมื่อค่า $t_2$ เข้าใกล้ค่า $t_1$ เรื่อย ๆ \pa เราจะได้ความเร็วของวัตถุในเวลา $t_1$ เป็น \pa
\[ \lim_{t_2 \rightarrow t_1} \frac{s(t_2) - s(t_1)}{t_2 - t_1} = s\,'(t_1) \] \pa

เราเขียนแทนด้วยด้วย $v(t) = s\,'(t)$
\end{slide}
\end{frame}

\begin{frame}[t]
\frametitle{ความสัมพันธ์ระหว่างความเร็ว $v$ และความเร่ง $a$}
ความเร่งเฉลี่ยของวัตถุจาก $t = t_1$ ถึง $t = t_2$ คือ\pa 
\begin{slide}
\[ \frac{v(t_2) - v(t_1)}{t_2 - t_1}\] \pa

ถ้าเราหาลิมิตเมื่อค่า $t_2$ เข้าใกล้ค่า $t_1$ เรื่อย ๆ \pa เราจะได้ความเร่งของวัตถุในเวลา $t_1$ เป็น \pa
\[ \lim_{t_2 \rightarrow t_1} \frac{v(t_2) - v(t_1)}{t_2 - t_1} = v\,'(t_1)\] \pa

เราเขียนแทนด้วยด้วย $a(t) = v\,'(t)$ \pa

สังเกตว่า $a(t) = \frac{d}{dt} v(t) = \frac{d^2}{dt^2} s(t)$ ซึ่งเป็นอนุพันธ์อันดับสองของ $s(t)$
\end{slide}
\end{frame}

\begin{frame}
\begin{thm}{}{}
กำหนดให้อนุภาคเคลื่อนที่ตามแนวเส้นตรง โดยที่ 
\begin{itemize}
\item $s(t)$ แทนตำแหน่งของอนุภาค ณ เวลา $t$
\item $v(t)$ แทนความเร็วของอนุภาค ณ เวลา $t$
\item $a(t)$ แทนความเร่งของอนุภาค ณ เวลา $t$
\end{itemize}
แล้วจะได้ความสัมพันธ์ดังนี้
\begin{eq}
v(t) &=& s'(t) = \frac{d}{dt} s(t) \label{eq:applied_velo}\\
a(t) &=& v'(t) = \frac{d}{dt} v(t) \label{eq:applied_accel1}\\
&=& s''(t) = \frac{d^2}{dt^2} s(t) \label{eq:applied_accel2}
\end{eq}
\end{thm}
\end{frame}


\begin{frame}[t]
\begin{ex}
ถ้ารถวิ่งเป็นเส้นตรง โดยมีสมการของการกระจัดเป็น $s(t) = t^2+5t$ ซึ่งมีหน่วยเป็นเมตร และเวลา $t$ มีหน่วยเป็นนาที จงหาความเร็วและความเร่งของรถคันนี้ที่เวลา $t= 3$
\end{ex}
\begin{sol}
จากสมการ \eqref{eq:applied_velo} และ \eqref{eq:applied_accel1}จะได้ว่า
\begin{eq*}
v(t) &=& s'(t) = \frac{d}{dt} (t^2+5t) = 2t+5 \\
a(t) &=& v'(t) = \frac{d}{dt} (2t+5) = 2
\end{eq*}
เมื่อแทนค่า $t=3$ จะได้
\begin{eq*}
v(3) &=& 2(3) + 5 = 11 \\
a(3) &=& 2
\end{eq*}
นั่นคือ ความเร็วเท่ากับ $3$ เมตรต่อนาที และความเร่งเท่ากับ $2$ เมตรต่อนาที
\end{sol}
\end{frame}


\begin{frame}[t]
\begin{ex}
ลูกบอลเหล็กติดปลายสปริง เคลื่อนที่โดยมีสมการของการกระจัดเป็น $s(t) = \sin (3t)$ ซึ่งมีหน่วยเป็นเซนติเมตร และเวลา $t$ มีหน่วยเป็นวินาที จงหาความเร็วและความเร่งของลูกบอลเหล็ก ณ เวลา $t = \pi$
\end{ex}
\begin{sol}
จากสมการ \eqref{eq:applied_velo} และ \eqref{eq:applied_accel1}จะได้ว่า
\begin{eq*}
v(t) &=& s'(t) = \frac{d}{dt} (\sin (3t)) = 3\cos (3t) \\
a(t) &=& v'(t) = \frac{d}{dt} (3\cos (3t)) = -9\sin(3t)
\end{eq*}
เมื่อแทนค่า $t=\pi$ จะได้
\[v(\pi) = 3\cos(3 \pi) = -3 \qquad a(\pi) = -9\sin(3 \pi) = 0\]
นั่นคือ ลูกบอลเหล็กมีความเร็วเท่ากับ $-3$ เซนติเมตรต่อนาที และความเร่งเท่ากับ $0$ เซนติเมตรต่อนาทีกำลังสอง โดยความเร็วที่ติดลบหมายความว่า ณ เวลานั้น ลูกบอลเหล็กกำลังเคลื่อนที่ในทิศทางลบ
\end{sol}
\end{frame}



\begin{frame}[t]
\begin{ex}
ถ้าโยนลูกบอลขึ้นฟ้า โดยมีสมการของการกระจัดเป็น $s(t) = 20t-5t^2$ ซึ่งมีหน่วยเป็นเมตร และ $t$ มีหน่วยเป็นวินาที จงหาว่า
\begin{tasks}(1)
\task ลูกบอลจะกลับมาที่จุดเริ่มต้นเมื่อเวลาผ่านไปเท่าใด
\task ณ เวลาที่ลูกบอลกลับมาที่จุดเดิม ลูกบอลจะมีความเร็วและความเร่งเท่าใด
\end{tasks}
\end{ex}
\end{frame}

\begin{frame}
\begin{sol}
\begin{tasks}(1)
\task ลูกบอลจะกลับมาที่จุดเริ่มต้นเมื่อ $s(t)=0$ นั่นคือเราแก้สมการ $20t-5t^2=0$ นั่นคือ $t = 0,4$ แต่ $t=0$ คือเวลาที่ลูกบอลถูกโยนออก ฉะนั้นลูกบอลจะกลับมาที่จุดเดิมหลังผ่านไป 4 วินาที
\task จากสมการ \eqref{eq:applied_velo} และ \eqref{eq:applied_accel1}จะได้ว่า
\begin{eq*}
v(t) &=& s'(t) = \frac{d}{dt} (20t-5t^2) = 20-10t \\
a(t) &=& v'(t) = \frac{d}{dt} (20-10t) = -10
\end{eq*}
เมื่อแทนค่า $t=4$ ซึ่งเป็นเวลาที่ลูกบอลกลับมาที่จุดเริ่มต้น จะได้ $v(4) = 20-10(4) = -20$ และ $a(4) = -10$ นั่นคือ ลูกบอลมีความเร็วเท่ากับ $-20$ เมตรต่อวินาที และความเร่งเท่ากับ $-10$ เมตรต่อนาทีกำลังสอง
\end{tasks}
\end{sol}
\end{frame}

\begin{frame}
\frametitle{หมายเหตุ}
เราใช้อนุพันธ์เพื่ออธิบายความสัมพันธ์ได้ดังนี้
\[ s(t) \xrightarrow{\frac{d}{dt}} v(t) \xrightarrow{\frac{d}{dt}} a(t)\]
นั่นคือ ถ้าเราทราบ $s(t)$ แล้วเราจะสามารถหา $v(t)$ และ $a(t)$ ได้ แต่ในทางกลับกัน ถ้าให้ $a(t)$ มาให้แล้วเราจะย้อนกลับไปหา $s(t)$ ได้หรือไม่
\[ s(t) \xleftarrow{?} v(t) \xleftarrow{?} a(t)\]
กระบวนการย้อนกลับนี้ เราเรียกว่า ``ปริพันธ์'' และจะได้ศึกษาในหน่วยต่อไป
\end{frame}

%%%%%%%%%%%%%%
\section{อัตราสัมพัทธ์}
\frame{\sectionpage}
%%%%%%%%%%%%%%

\begin{frame}
\frametitle{ที่มา}
อนุพันธ์คือการหาอัตราความเปลี่ยนแปลงของตัวแปร 2 ตัว \pa
\[y = f(x)\]
ในโจทย์บางประเภท จะมีตัวแปรที่สามเพิ่มขึ้นมา เช่น
\begin{tasks}[style=enumerate](1)
\task ราคาสินค้าขึ้นอยู่กับคุณภาพของสินค้า และคุณภาพของสินค้าขึ้นอยู่กับต้นทุน
\task ความเร็วของรถขึ้นอยู่กับแรงม้าของเครื่องยนต์ และแรงม้าของเครื่องยนต์ขึ้นอยู่กับราคา
\task น้ำหนักของวัสดุขึ้นอยู่กับปริมาตร และปริมาตรของวัตถุขึ้นอยู่กับอุณหภูมิ
\end{tasks}
\end{frame}

\begin{frame}[t]
สมมติว่าตัวแปร $x,y$ และ $u$ มีความสัมพันธ์ดังนี้
\[y = f(u), \qquad u = g(x)\]
\begin{slide}
ซึ่งก็คือฟังก์ชันประกอบนั่น
\[y = f(g(x))\]
ฉะนั้นเราสามารถหาอัตราสัมพัทธ์ $y$ เทียบกับ $x$ ได้ด้วยกฎลูกโซ่
\[ \dyx = \dx f(g(x)) = \left( \frac{d}{du} f(u) \right) \left( \frac{d}{dx} g(x) \right)\]
โจทย์อัตราสัมพัทธ์ จะให้ค่าของ $\frac{dy}{dx}, \frac{d}{du} f(u)$ หรือ $\frac{d}{dx} g(x)$ มาค่าใดค่าหนึ่ง ซึ่งเราสามารถหาค่าที่สองได้ และโจทย์จะถามถึงค่าที่สาม โดยดูได้จากตัวอย่างดังต่อไปนี้
\end{slide}
\end{frame}


\begin{frame}[t]
\begin{ex}
สี่เหลี่ยมจัตุรัสรูปหนึ่ง มีอัตราการขยายด้านทุกด้านเป็น 3 เซนติเมตรต่อวินาที จงหาอัตราการเปลี่ยนแปลงของพื้นที่ต่อเวลา
\end{ex}
\begin{sol}
จากโจทย์เรากำหนดตัวแปรได้ดังนี้
\begin{tasks}[style=enumerate](1)
\task $a$ เป็นพื้นที่ของสี่เหลี่ยมจัตุรัส 
\task $s$ เป็นความยาวด้านของสี่เหลี่ยมจัตุรัส
\task $t$ เป็นเวลา
\end{tasks}
โจทย์กำหนดอัตราการขยายของด้าน ซึ่งเขียนโดยใช้ตัวแปรที่เรากำหนดได้เป็น $\frac{ds}{dt} = 3$ ซึ่งโจทย์ต้องการหา $\frac{da}{dt}$ โดยกฎลูกโซ่จะได้
\[ \frac{da}{dt} = \left( \frac{da}{ds} \right) \left( \frac{ds}{dt} \right)\]
\end{sol}
\end{frame}

\begin{frame}
\begin{sol} (ต่อ)
นั่นคือ เราต้องการหา $\frac{da}{ds}$ ซึ่งโจทย์ไม่ได้กำหนดมาให้ แต่เราทราบความสัมพันธ์ของพื้นที่และความยาวด้านของสี่เหลี่ยมจัตุรัสว่า
\[a = s^2\]
ฉะนั้น
\[ \frac{da}{ds} = \frac{d}{ds} (s^2) = 2s\]
นั่นคือ
\[ \frac{da}{dt} = \left( \frac{da}{ds} \right) \left( \frac{ds}{dt} \right) = (2s)(3) = 6s\]
สังเกตว่า แม้ว่าจะเป็นอัตราสัมพัทธ์ของพื้นที่และเวลา แต่มีตัวแปรเป็นความยาวด้าน
\end{sol}
\end{frame}


\begin{frame}[t]
\begin{ex}
ชายผู้หนึ่งเดินทางเข้าหาฐานหาคอยสูง 60 เมตร ด้วยความเร็ว 2 เมตรต่อวินาที จงหาว่าชายผู้นี้จะเดินเข้าใกล้ยอดหอคอยด้วยอัตราเร็วเท่าใด ขณะที่เขาอยู่ห่างจากฐานหอคอยเป็นระยะ 8 เมตร
\end{ex}
\begin{sol}
จากโจทย์เรากำหนดตัวแปรได้ดังนี้
\begin{tasks}[style=enumerate](1)
\task $r$ เป็นระยะทางระหว่างชายและยอดหอคอย 
\task $s$ เป็นระยะทางระหว่างชายและฐานหอคอย
\task $t$ เป็นเวลา
\end{tasks}
ซึ่งโจทย์กำหนดให้ $\frac{ds}{dt} = 2$ และอยากทราบ $\frac{dr}{dt}$ ซึ่งจากกฎลูกโซ่จะได้
\[\frac{dr}{dt} = \left(\frac{dr}{ds}\right) \left(\frac{ds}{dt}\right)\]
\end{sol}
\end{frame}

\begin{frame}
\begin{sol} (ต่อ)
แม้โจทย์จะไม่กำหนดความสัมพันธ์ระหว่าง $r$ และ $s$ มาให้ แต่เราทราบจากทฤษฎีบทปีทากอรัสว่า 
\[r^2 = s^2 + 60^2\] 
นั่นคือ 
\[r = \sqrt{ s^2 + 3600} = (s^2+3600)^{1/2}\]
ฉะนั้น อัตราสัมพัทธ์ของ $r$ เทียบกับ $s$ คือ
\begin{eq*}
\frac{dr}{ds} &=& \frac{d}{dx} (s^2 + 3600)^{1/2} \\
&=& \frac{1}{2}(s^2+3600)^{-1/2} \frac{d}{ds} (s^2+3600) \mathnote{กฎลูกโซ่ โดยกำหนดให้ $u = s^2+3600$}\\
&=& \frac{2s}{2\sqrt{s^2+3600}} = \frac{s}{\sqrt{s^2+3600}}
\end{eq*}
\end{sol}
\end{frame}

\begin{frame}
\begin{sol} (ต่อ)
ดังนั้น
\[\frac{dr}{dt} = \left(\frac{dr}{ds}\right) \left(\frac{ds}{dt}\right) = \left( \frac{s}{\sqrt{s^2 + 3600}} \right) (2) = \frac{2s}{\sqrt{s^2 + 3600}}\]
เมื่อเขาอยู่ห่างจากฐานหอคอยเป็นระยะ 8 เมตร นั่นคือเมื่อ $s = 8$ จะได้อัตราสัมพัทธ์เป็น
\[ \frac{2(8)}{\sqrt{8^2 + 3600}} = \frac{16}{\sqrt{3664}} \approx 0.0043 \]
\end{sol}
\end{frame}


\begin{frame}[t]
\begin{ex}
ถ้าสูบลมเข้าไปในลูกบอลด้วยอัตราเร็วคงที่ 15 ลูกบาศก์เซนติเมตรต่อวินาที สมมติให้อากาศมีความดันคงที่ จงหาอัตราการเปลี่ยนแปลงของรัศมีที่เพิ่มขึ้นเมื่อลูกบอลมีรัศมี 10 เซนติเมตร
\end{ex}
\begin{sol}
จากโจทย์เรากำหนดตัวแปรได้ดังนี้ $v$ เป็นปริมาตรของลูกบอล, $r$ เป็นรัศมีของลูกบอล และ $t$ เป็นเวลา

ซึ่งโจทย์กำหนดให้ $\tfrac{dv}{dt} = 15$ และอยากทราบ $\tfrac{dr}{dt}$ ซึ่งจากกฎลูกโซ่จะได้
\[\frac{dv}{dt} = \left(\frac{dv}{dr}\right) \left(\frac{dr}{dt}\right)\]
สังเกตว่า ในโจทย์ข้อนี้เรามีข้อมูลฝั่งซ้ายของสมการ และต้องการหาค่าของ $\frac{dr}{dt}$ ซึ่งอยู่ฝั่งขวาของสมการแทน
\end{sol}
\end{frame}

\begin{frame}
\begin{sol} (ต่อ) เราทราบความสัมพันธ์ระหว่างรัศมีและปริมาตรเป็น
\[v = \tfrac{4}{3}\pi r^3\]
ฉะนั้น เมื่อหาอนุพันธ์จะได้
\[\tfrac{dv}{dr} = \tfrac{d}{dr} \left( \tfrac{4}{3}\pi r^3 \right) = 4\pi r^2\]
ฉะนั้น
\begin{eq*}
\frac{dv}{dt} &=& \left(\frac{dv}{dr}\right) \left(\frac{dr}{dt}\right) \\
15 &=& (4\pi r^2) \left( \frac{dr}{dt} \right) \\
\frac{dr}{dt} &=& \frac{15}{2\pi r^2}
\end{eq*}
เมื่อแทนค่า $r =10$ จะได้ 
\[\tfrac{dr}{dt} = \tfrac{15}{2 \pi (10)^2} = \tfrac{3}{40 \pi} \approx 0.012\]
\end{sol}
\end{frame}


%%%%%%%%%%%%%%
\section{การหาค่าประมาณ}
\frame{\sectionpage}
%%%%%%%%%%%%%%

\begin{frame}[t]
\frametitle{ที่มา}
จากนิยามอนุพันธ์
\[ f\,'(x) = \limh \frac{f(x+h) - f(x)}{h}\] \pa
\begin{slide}
แปลว่าถ้าค่า $h$ เข้าใกล้ 0 มาก ๆ \pa ค่า $\frac{f(x+h)-f(x)}{h}$ ก็จะเข้าใกล้ค่า $f\,'(x)$ 

กำหนดให้ $\Delta x$ เป็นความเปลี่ยนแปลง ซึ่งเป็นค่าที่น้อยมาก ๆ นั่นคือ เราสามารถบอกได้ว่า \pa
\[ f\,'(x) \approx \frac{f(x+\Delta x) - f(x)}{\Delta x}\pa\] \pa

หรือจัดรูปได้เป็น
\begin{eq*}
f\,'(x)\Delta x &\approx& f(x+\Delta x) - f(x) \\
f(x+\Delta x) &\approx& f\,'(x) \Delta x + f(x)
\end{eq*}
\end{slide}
\end{frame}



\begin{frame}
\frametitle{ค่าเชิงอนุพันธ์}
\begin{defn}{}{}
ให้  $y= f(x)$ เป็นฟังก์ชันที่หาอนุพันธ์ได้ที่ $x$  

เรียกพจน์ $f\,'(x) \Delta x$ ว่า \textbf{ค่าเชิงอนุพันธ์} หรือ \textbf{ผลต่างอนุพันธ์}ของ $f$  ที่ $x$  และเขียนแทนด้วยสัญลักษณ์ $df$ หรือ $dx$
\end{defn} \pa
\end{frame}

\begin{frame}
ในคาบเรียนก่อนเราได้ศึกษาถึงสมการเส้นสัมผัสเส้นโค้ง เราจะนำสมการนั้นมาใช้อีกครั้งหนึ่งที่นี่
\begin{defn}{}{}
กำหนดฟังก์ชัน $f(x)$ ที่่หาอนุพันธ์ได้ที่จุด $x_0$ นิยาม\textbf{ฟังก์ชันค่าประมาณเชิงเส้น}เป็น
\begin{equation}
L(x) = f\,'(x_0) (x-x_0) + f(x_0)
\end{equation}
\end{defn}
สังเกตว่า $L(x_0) = f(x_0)$ นั่นคือสมการเส้นตรง $L$ ให้ค่าเท่ากับฟังก์ชัน $f$ ที่จุด $x_0$
\end{frame}

\begin{frame}
\addfig{0.8}{diff_applied_approx_intro}{การใช้เส้นตรง $L$ ประมาณค่าของฟังก์ชัน $f$}{}
\end{frame}

\begin{frame}
\begin{thm}{}{}
ถ้า $f(x)$ หาอนุพันธ์ได้ที่ $x$ และ $\Delta x$ เป็นค่าที่น้อยมากเมื่อเทียบกับ $x$ \pa แล้วจะได้
\[ f(x+\Delta x) \approx f(x) + f\,'(x) \Delta x\]
\end{thm}\pa
ทฤษฎีบทนี้มีประโยชน์ในการประมาณค่าของฟังก์ชัน ที่ใกล้เคียงกับค่าที่เรารู้อยู่แล้ว \pa นั่นคือ เราใช้ค่า $f(x)$ มาประมาณค่า $f(x+\Delta x)$ โดยมีผลต่างโดยประมาณเป็น $f\,'(x) \Delta x$
\end{frame}
 

\begin{frame}[t]
\begin{ex}
จงประมาณประมาณค่า $20.01^3$ โดยใช้ค่าเชิงอนุพันธ์
\end{ex}\pa
\begin{sol}
กำหนดให้ $f(x) = x^3$ \pa โดยเลือก $x = 20$ และ $\Delta x = 0.01$ \pa ทำให้ 
\[20.01^3 = (20 + 0.01)^3 \pa = f(20+0.01) = f(x + \Delta x)\]
จะได้ว่า 
\begin{eqnarray*}
f(x+\Delta x) &\approx& f(x) + f\,'(x) \Delta x \pa \\
&\approx& f(20) + f\,'(20) (0.01) \pa \\
&\approx& 8000 + 3(20)^2 (0.01) \pa \\
&=& 8000 + 12 = 8012  \pa
\end{eqnarray*}
นั่นคือ $20.01^3 \approx 8012$ \pa ซึ่งเครื่องคิดเลขบอกเราว่า $20.01^3 = 8012.006001$
\end{sol}
\end{frame}


\begin{frame}[t]
\begin{ex}
จงประมาณประมาณค่า  $\sqrt{4.2}$ โดยใช้ค่าเชิงอนุพันธ์
\end{ex}\pa
\begin{sol}
กำหนดให้ $f(x) = \sqrt{x}$ \pa โดยเลือก $x = 4$ และ $\Delta x = 0.2$ \pa ทำให้ 
\[\sqrt{4.2} = f(4+0.2) = f(x + \Delta x)\] \pa
จะได้ว่า 
\begin{eqnarray*}
f(x+\Delta x) &\approx& f(x) + f\,'(x) \Delta x \pa \\
&\approx& f(4) + f\,'(4) (0.2) \pa \\
&\approx& 2 + \frac{1}{2\sqrt{4}} (0.2) = 2.05 \pa
\end{eqnarray*}
นั่นคือ $\sqrt{4.2} \approx 2.05$ \pa ซึ่งเครื่องคิดเลขบอกเราว่า $\sqrt{4.2} \approx 2.04939$
\end{sol}
\end{frame}


\begin{frame}[t]
\begin{ex}
กำหนดฟังก์ชัน $f(x) = e^x$ จงหาค่าประมาณของ $f(h)$ เมื่อ $h$ มีค่าใกล้ $0$ โดยใช้ค่าเชิงอนุพันธ์
\end{ex}\pa
\begin{sol}
กำหนดให้ $f(x) = e^x $ \pa โดยเลือก $x = 0$ และ $\Delta x = h$ \pa จะได้
\[ f\,'(x) = \dx e^x = e^x  \rightarrow  f\,'(0) = e^0 = 1\]
ฉะนั้นว่า 
\begin{eqnarray*}
f(x+\Delta x) &\approx& f(x) + f\,'(x) \Delta x \pa \\
f(0+h)&\approx& f(0) + f\,'(0) (h) \pa \\
e^h &\approx& 1 + h
\end{eqnarray*}
นั่นคือ เมื่อ $h$ มีค่าใกล้ 0 แล้ว $f(h) \approx 1+h$
\end{sol}
\end{frame}